\textbf{Parameter-efficient fine-tuning.} Houlsby et al.~\cite{houlsby2019parameter} insert small bottleneck adapters into a frozen transformer; BitFit~\cite{benzaken2021bitfit} trains only bias terms; LoRA~\cite{hu2021lora} adds trainable low-rank factors to weight matrices and adds no inference latency after merging. He et al.~\cite{he2021unified} analyse these methods as modifications of hidden states and study design choices that transfer between them. DoRA~\cite{liu2024dora} starts from the observation of an accuracy gap between LoRA and full fine-tuning, decomposes weights into magnitude and direction, and applies LoRA to the direction. Our last-block and head-only baselines are the simplest ``train a few layers'' alternatives and are not taken from these papers.

\textbf{Why low rank.} Aghajanyan et al.~\cite{aghajanyan2020intrinsic} show that pretrained language models can be fine-tuned in a random subspace of very low dimension, which is the usual motivation for low-rank updates. Xu et al.~\cite{xu2025understanding} analyse LoRA dynamics for matrix factorisation under gradient flow and show that the initialisation scale affects the final error; this is relevant to our observation that optimisation of low-rank factors can be fragile, though our setting is not matrix factorisation and we do not test their account.

\textbf{Forgetting.} Catastrophic forgetting has long been studied in continual learning; EWC~\cite{kirkpatrick2016overcoming} penalises changes to weights important for earlier tasks. Biderman et al.~\cite{biderman2024lora} compare LoRA and full fine-tuning on programming and mathematics for instruction fine-tuning and continued pretraining and report that LoRA substantially underperforms full fine-tuning in standard low-rank settings but better maintains performance outside the target domain. Our study is the closest synthetic analogue we know of; it differs in scale (70k parameters, 2000 examples), in measuring forgetting of a single known task, and in giving the model a task token.

\textbf{Optimisation.} We train with AdamW~\cite{loshchilov2017decoupled} using weight decay zero. We implement every method here and reuse no released code.
